\section{Onslaught at Bellosguardo}

It was three o’clock in the morning of November first. One hundred and fifty delighted guests had left my masked ball. It was a ritual I hosted every 31st of October in the Great Hall of my Villa of the Saracen, an architectural jewel on the Piazza di Bellosguardo, Florence.

``My dear, I thought your idea of playing Saint Saens’ ``Danse Macabre'' on your piano at the stroke of midnight was great theater,'' declared Sir Harold Acton, Lord and Master of the 16th century Florentine estate Villa La Pietra.

Father Andreas Resch, the psychiatrist, exorcist and demonologist, bent down to give me a hug. Andreas was my mentor, confessor and friend.

``You are spiritually vulnerable, you are aware of that yes? Ring me at the convent if anything strange occurs,'' he murmured.

``Really Andreas, the creatures and I have established a tolerant modus vivendi.''

``Don’t be flippant about this Isabella. Promise me you’ll be careful.''

The granite fireplace in my bedroom reached up to the vaulted ceiling. It cast suggestive shadows on the walls, which fascinated me. My bedroom, music room and library stood at the far end of the Great Hall.

I shall undress by the fire. The air has suddenly turned gelid.

In 1519, the Bardi banking family, deadly rivals of the Medici, commissioned Baccio D’Agnolo, master and mentor of young Michelangelo to design a villa in the Florentine hills.

``It must surpass all the Medici dwellings in Florence,'' commanded the Bardi. The Villa of the Saracen not only had Florence at her feet, but all her splendors proudly stood facing her.

More than four centuries have passed.  I can’t believe she’s mine.

``Contessa, you are willful and capricious. The villa has naughty and nasty spirits who drive everyone away.''  I was sternly warned.

``I don’t care. I’ll charm and seduce them with music and dance. I’ll surround it with my beauty and laughter. Nothing is going to stop me from living in that enchanted villa,'' I replied resolutely.

The ``beings'' and I established an affectionate entente. They opened and shut massive wood drawers in the Great Hall only between the hours of 9:00 and 10:00 in the evening. The velvet curtains in the Dining Hall swished to and fro at dusk and stopped abruptly the moment I said, ``guys and dolls, enough. You have broken my chops.'' The howling wind came calling when I was cozily inside my quilted silk embroidered coverlet. Faces without bodies and vice versa swirled in the vaulted ceilings. One of them resembled Franz Liszt, my favorite composer. He had spent a summer in Bellosguardo. It was natural for him to fly around, particularly when I played his sublime but taxing Transcendental Studies.

I stepped out of my opulent scarlet gown and inched closer to the fire. I stood there, naked. Its warmth filled me with a sensual torpor. The mewing sounds coming from the Great Hall startled me.

I forgot to close the door. What are my cats doing out there? Oh! They’re here, sleeping on my bed. Someone is still crying, one of the children? Criminy!  Here I am illuminating my nakedness.

``Marko, Cynthia, what’s wrong? I’ll be right there.'' I called out, hastily grabbing my nightgown where it lay across the bed and slipping it over my head. I stood on the threshold between my bedroom and the Great Hall. A viscid fog reeking of sweat and sexual fluids assaulted me. Invisible hands lifted my long nightgown soaked with coital frenzy. A force was pushing me slowly out of my bedroom.

Mistake. I had addressed entities indulging in the most wanton sex. I heard the moans of a woman, accompanied by the loud rhythmic creaking of a high backed carved chair. A man gasped raucously. Every sound reverberated throughout the vaulted ceilings of the Great Hall.

``No! No!'' I yelled in horror and ran towards the telephone by my bed. Oh dear Jesus. I can’t remember Andreas’s telephone number. Isabella, be calm. Recite the Our Father. Don’t listen to the sounds of lust. Ah, I got the number right at last. It’s ringing.

``Cara, I was hoping it wasn’t you.''

``Andreas, I’m under attack. I have never felt such terror. I fear for my soul. Help me please.''

``Isabella, listen to me. Do not venture into the Great Hall under any circumstances. Surround yourself with light. You can overcome them with love and prayer.

``How long will this onslaught last Andreas?''

``I don’t know. My prayers and rituals will help you. Don’t be afraid, whatever happens. Trust in Jesus.''

I turned on all three chandeliers and lit every candle. I placed seven blue candles on my night table beside the crucifix — for the Seven Last Words of Jesus on the cross. I knelt on the Florentine tiles and repeated them over and over as loudly as I could without damaging my vocal chords.

In the Great Hall the lovers continued their crazed sex. They screamed and yelled in lust.

I remained close to the light and the crucifix. I dared not close the door and expose myself to the forces of darkness.

The sobs and sighs abated and finally vanished.

Deliverance!

``There are three possible explanations for what you experienced in the darkest hours of dawn. You entered another dimension that occurred centuries ago and received their signals much like a radio. Because of your long sexual abstinence, your powerful psyche might have created the frenetic coupling. Demonic entities tried to possess you by using a weapon that almost always works — sex.''

``I thought I knew most spirits at the Villa.''

``After 450 years, I doubt you’ll ever encounter all of them but never let your guard down. Is the chapel ready to receive the Blessed Sacrament?''

``Yes. My children and I will assist you during the Mass on this special All Saints Day.''

``Sia Lodato Jesu Cristo.''

``Amen,'' I humbly replied.

THE END

N. B. This is a true story in all its details. Real names and places have been used throughout the narrative. Contessa Isabella von Vacani.